\scanpage{181}
est l'hom
\[
\begin{aligned}
f_*(\check E^{(p)}\otimes\Omega^n_{X/S})&\longrightarrow
 g_*(\check{\widetilde E}\otimes\Omega^n_{X^{p/S}/S})\\
\deleted{f_*(\check E\otimes\Omega}&
 \quad\rotatebox{90}{$\simeq$}\uparrow\ \text{\unsure{ss|sous} réserve de commutation à $F_S$ de $f_*(\check E\otimes\Omega^n_{X/S})$}\\
&g_*(\check E\otimes\Omega^n_{X/S})^{(p)}
\end{aligned}
\]
donné par l'opération de Cartier $(\ast\ast)$ ci-dessus.

Le cas le plus important est celui où $E=\underline O_X$, on trouve que l'hom de puissance p.ième
\[
 R^nf_*(\underline O_X)^{(p)}\mathrel{\text{---}}R^nf_*(\underline O_X)
\]
est transposée de l'opération de Cartier
\[
 \widetilde C:f_*(\Omega^n_{X/S})\mathrel{\text{---}}g_*(\Omega^n_{X^{p/S}/S})
 \mathrel{\underset{\text{ss. réserve}}{\overset{\sim}{\longleftarrow}}}
 f_*(\Omega^n_{X/S})^{\otimes p}
\]
$\overline{\mathrm M}$ énoncé pour les $R^if_*(\underline O_{X\deleted{\unsure{{}^{p/S}}}})^{(p)}\mathrel{\text{---}}R^if_*(\underline O_X)$, \inserted{sous réserve que c'est $R^if_*(\underline O_{X^{p/S}/S})$} à cela près qu'il faut prendre l'hom sur $R^{n-\unsure{1|i}}f_*(\Omega^n_{X/S})\longrightarrow R^{n-i}f_*(\Omega^n_{X/S})^{\unsure{\otimes p|\otimes1/p|(p)}}$ \inserted{[sous réserve]} déduit de $\widetilde C$ par application de $R^{n-i}f_*$.

Un hom intéressant est
\[
 R^1f_*(\underline O_{X^{p/S}})\longrightarrow R^1f_*(\underline O_X)
\]
dans le cas où $f$ est propre, \deleted{car} \deleted{\unsure{c'est}} il s'identifie à l'hom. des sections pour les fibrés vectoriels
\[
 \underset{\substack{\parallel\\(\underline{\operatorname{Lie}}\,\underline{\operatorname{Pic}}_{X/S})^{(p/S)}}}
 {\underline{\operatorname{Lie}}\,\underline{\operatorname{Pic}}_{X^{p/S}/S}}
 \longrightarrow\underline{\operatorname{Lie}}\,\underline{\operatorname{Pic}}_{X/S}
\]

\margin{\uline{À vérifier} (et y donner un sens sous des conditions aussi générales que possible)}
% slanted
Nous admettrons ici que \deleted{\unsure{l'}} l'hom \deleted{\unsure{sur Pic}}

\diagram{\begin{tikzcd}[ampersand replacement=\&]
\underline{\operatorname{Pic}}_{X^{p/S}/S}\arrow[r,no head]\arrow[d,no head,"\rotatebox{90}{$\simeq$}" description,draw=none]\&\underline{\operatorname{Pic}}_{X/S}\\
\underline{\operatorname{Pic}}_{X/S}^{p/S}\arrow[ur,"V"]\&
\end{tikzcd}}
\margin{Il est immédiat en tout cas que $F_{\mathrm{Pic}}V=p\,\mathrm{id}$, $VF_{\mathrm{Pic}}=p.\mathrm{id}$ \deleted{\unsure{ce qui}}}
est le \uline{Verschiebung} , d'où il résulte que l'hom.
% archive sheet number: 180
